\section{Ход работы}

\subsection{Таблица истинности (пункт~\ref{task:table})}
Число $A = 2a_1 + a_0$ принимает значения от 0 до 3, как и $B$, поэтому у функции 16
наборов входов. Функция равна 1, когда $A > B$. Таблица истинности приведена в
таблице~\ref{tab:truth}: единиц в ней шесть, по числу пар $A > B$.

\begin{bbttable}{Таблица истинности функции $F = (A > B)$}{tab:truth}
  \begin{tabular}{|c|c|c|c|c|c|c|}
    \hline
    $a_1$ & $a_0$ & $b_1$ & $b_0$ & $A$ & $B$ & $F$ \\ \hline
    0 & 0 & 0 & 0 & 0 & 0 & 0 \\ \hline
    0 & 0 & 0 & 1 & 0 & 1 & 0 \\ \hline
    0 & 0 & 1 & 0 & 0 & 2 & 0 \\ \hline
    0 & 0 & 1 & 1 & 0 & 3 & 0 \\ \hline
    0 & 1 & 0 & 0 & 1 & 0 & 1 \\ \hline
    0 & 1 & 0 & 1 & 1 & 1 & 0 \\ \hline
    0 & 1 & 1 & 0 & 1 & 2 & 0 \\ \hline
    0 & 1 & 1 & 1 & 1 & 3 & 0 \\ \hline
    1 & 0 & 0 & 0 & 2 & 0 & 1 \\ \hline
    1 & 0 & 0 & 1 & 2 & 1 & 1 \\ \hline
    1 & 0 & 1 & 0 & 2 & 2 & 0 \\ \hline
    1 & 0 & 1 & 1 & 2 & 3 & 0 \\ \hline
    1 & 1 & 0 & 0 & 3 & 0 & 1 \\ \hline
    1 & 1 & 0 & 1 & 3 & 1 & 1 \\ \hline
    1 & 1 & 1 & 0 & 3 & 2 & 1 \\ \hline
    1 & 1 & 1 & 1 & 3 & 3 & 0 \\ \hline
  \end{tabular}
\end{bbttable}

\subsection{СДНФ (пункт~\ref{task:sdnf})}
По шести единицам таблицы~\ref{tab:truth} СДНФ записывается формулой~\eqref{eq:sdnf}:
\begin{multline}\label{eq:sdnf}
  F = \overline{a_1}\, a_0\, \overline{b_1}\, \overline{b_0}
  \vee a_1\, \overline{a_0}\, \overline{b_1}\, \overline{b_0}
  \vee a_1\, \overline{a_0}\, \overline{b_1}\, b_0 \vee {} \\
  {} \vee a_1\, a_0\, \overline{b_1}\, \overline{b_0}
  \vee a_1\, a_0\, \overline{b_1}\, b_0
  \vee a_1\, a_0\, b_1\, \overline{b_0}.
\end{multline}

Схема прямо по СДНФ потребовала бы шести элементов И на четыре входа и элемента ИЛИ на шесть
входов, поэтому функцию нужно минимизировать.

\subsection{Минимизация по карте Карно (пункт~\ref{task:min})}
Карта Карно функции с контурами склейки показана на рисунке~\ref{fig:karnaugh}. Строки
подписаны значениями $a_1 a_0$, столбцы — $b_1 b_0$ в коде Грея.

\bbtScreenshot[0.75\linewidth]{materials/karnaugh.png}{Карта Карно функции $F$ с
контурами}{fig:karnaugh}

Единицы покрываются тремя контурами:
\begin{itemize}
  \item квадрат из четырёх клеток в строках $a_1 a_0 = 11, 10$ и столбцах $b_1 b_0 = 00, 01$ —
  $a_1\overline{b_1}$: старший разряд $A$ равен 1, а у $B$ — 0;
  \item пара клеток в столбце $b_1 b_0 = 00$, строки $a_1 a_0 = 01, 11$ —
  $a_0\overline{b_1}\,\overline{b_0}$: $B = 0$, а $A$ нечётно;
  \item пара клеток в строке $a_1 a_0 = 11$, столбцы $b_1 b_0 = 00, 10$ через край карты —
  $a_1 a_0 \overline{b_0}$: $A = 3$, а $B$ чётно.
\end{itemize}
Минимальная ДНФ — формула~\eqref{eq:mdnf}:
\begin{equation}\label{eq:mdnf}
  F = a_1\overline{b_1} \vee a_0\overline{b_1}\,\overline{b_0} \vee a_1 a_0 \overline{b_0}.
\end{equation}

Вместо шести конъюнкций по четыре переменных осталось три конъюнкции из двух и трёх
переменных. Контуры пересекаются в клетке $A = 3$, $B = 0$, это допустимо: каждая единица
покрыта хотя бы одним контуром, а нулей контуры не захватывают.

\subsection{Схема в Logisim (пункт~\ref{task:scheme})}
Схема по формуле~\eqref{eq:mdnf} показана на рисунке~\ref{fig:scheme}.

\bbtScreenshot{materials/scheme.png}{Схема сравнения $A > B$ по минимальной ДНФ}{fig:scheme}

Два элемента НЕ дают $\overline{b_1}$ и $\overline{b_0}$. Каждый из трёх элементов И
реализует одну конъюнкцию формулы~\eqref{eq:mdnf}: сверху вниз — $a_1\overline{b_1}$,
$a_0\overline{b_1}\,\overline{b_0}$ и $a_1 a_0 \overline{b_0}$. Элемент ИЛИ объединяет их в
выход $F$. Всего в схеме шесть элементов, и ни у одного нет больше трёх входов.

\subsection{Проверка тестовыми векторами (пункт~\ref{task:test})}
Для проверки составлен файл тестовых векторов из листинга~\ref{lst:vectors}: по строке на
каждый набор входов с ожидаемым значением $F$ из таблицы~\ref{tab:truth}.

\bbtCodeFile[text]{materials/test-vectors.txt}{Тестовые векторы}{lst:vectors}

Результат проверки в Logisim показан на рисунке~\ref{fig:test}: все 16 наборов прошли, ошибок
нет. Значит, схема на рисунке~\ref{fig:scheme} реализует ту же функцию, что и таблица
истинности.

\bbtScreenshot[0.7\linewidth]{materials/test-result.png}{Результат проверки схемы тестовыми
векторами}{fig:test}
